\documentclass[../../../main.tex]{subfiles}
\begin{document}
In framework of many-body system, the fundamental object is the many-body wavefunction
\[
  \Psi(x_1,\ldots,x_N)
  =\Psi(\mathbf{r}_1,\sigma_1\dots,\mathbf{r}_N,\sigma_N)
\]
where the generalized coordinate $x = (\mathbf{r}, \sigma)$ not only takes into account the spatial coordinate $\mathbf{r}$ but also the intrinsic spin $\sigma$ pertaining to the particle.
The state of the system therefore lives in the tensor product Hilbert space $\mathcal{H}^{\otimes N}$.
For identical particles, the wavefunction must satisfy permutation symmetry,
\begin{equation*}
  \Psi(\ldots, x_i, \ldots, x_j, \ldots) = \pm \Psi(\ldots, x_j, \ldots, x_i, \ldots)
\end{equation*}

When the particle number becomes macroscopic $  N \rightarrow O(10^{23})$ and the system is described only through macroscopic observable rather than the full wavefunction, one replaces the pure state description $\Psi(x_1,\ldots,x_N)$ with an ensemble or density operator
\begin{equation*}
  \rho = \frac{e^{-\beta H}}{Z}, \quad Z = \text{Tr} \; e^{-\beta H}.
\end{equation*}
This quantum statistical mechanics step corresponds to assuming statistical equilibrium with a reservoir (canonical ensemble) or fixed energy (microcanonical ensemble).

Thermodynamics emerges when an additional limit is taken
\begin{equation*}
  N \to \infty, \quad V \to \infty, \quad \frac{N}{V} = \text{const.}
\end{equation*}
In this regime microscopic fluctuations become negligible compared with macroscopic averages.
Observables become smooth functions of intensive variables such as \(T, P,\) and \(\rho\).
Thermodynamic relations follow from derivatives of the free energy
\begin{equation*}
  F = -k_B T \ln Z, \quad U = -\frac{\partial}{\partial \beta} \ln Z
\end{equation*}
For example
\begin{equation*}
  S = -\left( \frac{\partial F}{\partial T} \right)_V, \quad P = -\left( \frac{\partial F}{\partial V} \right)_T
\end{equation*}
Thermodynamic quantities are obtained from the partition function,

Solid-state theory arises when the many-body system satisfies essential conditions
\begin{enumerate}
  \item The particles form a dense phase (typically electrons and ions).
  \item The ions form an approximately periodic lattice
        \begin{equation*}
          V(\mathbf{r} + \mathbf{R}) = V(\mathbf{r})
        \end{equation*}
        where \(\mathbf{R}\) is a Bravais lattice vector. In this condition, eigenstates satisfy Bloch's theorem $\psi_{nk}(\mathbf{r}) = e^{i \mathbf{k} \cdot \mathbf{r}} u_{nk}(\mathbf{r})$
  \item The electronic states can be treated in an effective single-particle framework due to screening and mean-field approximations.
\end{enumerate}

The many-body problem is then approximated by quasiparticles (electrons, phonons, magnons) moving in effective potentials.
Techniques such as Hartree–Fock, density functional theory, or Fermi liquid theory implement this reduction.

\subsection{One-particle Operator}
A one-particle operator  is an operator acting independently on each particle in a many-particle system.
For a system of $N$ particles with generalized coordinates $\mathbf{x}_i$ a general one-particle operator $\hat{O}_1$ has the form of sum over operators $\hat{o } ^{(1)}$, which acts on one particle
\begin{equation*}
  \hat{O}_1 = \sum_{i } \hat{o } ^{(1)}_{i}(\mathbf{x}_i)
\end{equation*}
For a system with many particles, this representation becomes inconvenient because every term must refer to a specific particle coordinate $\mathbf{x_i}$.
Second quantization introduce
\begin{equation*}
  \hat{O}_1 = \sum_{ij} \langle 1_i | \hat{o } ^{(1)} | 1_j \rangle \hat{a}_i^\dagger \hat{a}_j
\end{equation*}
with $\ket{\left\{ 1_i \right\}  }$ as one-particle basis whose creation and annihilation operators are related to other of the one-particle basis $\ket{\left\{ 1_j \right\} }$.
In the position basis, the matrix element can be written as
\begin{equation*}
  \langle 1_i|\hat{o } ^{(1)}|1_j\rangle = \int d^3 \mathbf{r} d^3 \mathbf{r}' u^*_{i}(\mathbf{r}) \langle \mathbf{r}|\hat{o } ^{(1)}|\mathbf{r}'\rangle u_j(\mathbf{r}')
\end{equation*}
As such, the representation in position basis reads 
\begin{equation*}
  \hat{O}_1 
  = \int d^3r\;d^3r'\; \psi^\dagger(\mathbf{r}) \langle \mathbf{r}|\hat{o}^{(1)}|\mathbf{r}'\rangle \psi(\mathbf{r}')
\end{equation*}

Assuming the operator $\hat{o} ^{(1)}$ is local
\begin{equation*}
  \langle \mathbf{r}|\hat{o}^{(1)}|\mathbf{r}'\rangle  = \hat{o } ^{(1)}(\mathbf{r },\mathbf{p})\delta(\mathbf{r }-\mathbf{r}')
\end{equation*}
the matrix elements turns diagonal
\begin{equation*}
  \langle 1_i|\hat{o } ^{(1)}|1_j\rangle = \int d^3 \mathbf{r} \;u^*_{i}(\mathbf{r}) \hat{o }^{1}(\mathbf{r},\mathbf{p}) u_j(\mathbf{r}')
\end{equation*}
and the representation is 
\begin{equation*}
  \hat{O}_1 = \int d^3r\;\psi^\dagger(\mathbf{r}) \langle \mathbf{r}|\hat{o}^{(1)}|\mathbf{r}\rangle \psi(\mathbf{r})
\end{equation*}

This representation has the following advantage.
\begin{enumerate}
  \item \textbf{Number particle independence.} The indices $(i,j)$ label single-particle basis states, computations are therefore performed in the space of orbitals. Consequently the complexity is controlled by the size of the one-particle basis rather than by the dimensionality of the $N$-particle wavefunction.
  \item \textbf{Straightforward basis transformation.} Changing the single-particle basis corresponds simply to a unitary transformation of the operators.
\end{enumerate}

\subsubsection{Example: particle density}
The particle density
\begin{equation*}
  \rho(\mathbf{r}) = \psi^\dagger(\mathbf{r})\psi(\mathbf{r}) = \sum_{\mathbf{k},\mathbf{k}'} \braket{\mathbf{k}|\mathbf{r} }\braket{\mathbf{r}|\mathbf{k}'} \hat{a}_\mathbf{k}^\dagger \hat{a}_{\mathbf{k'}}
\end{equation*}
has the Fourier transform of
\begin{equation*}
  \rho(\mathbf{q}) = \sum_{\mathbf{k}} \hat{a}_\mathbf{k}^\dagger \hat{a}_{\mathbf{k}+\mathbf{q}}= \sum_{\mathbf{k}} \hat{a}_{\mathbf{k}-\mathbf{q}/2}^\dagger \hat{a}_{\mathbf{k}+\mathbf{q}/2}
\end{equation*}
in plane wave basis. 
In Bloch state basis, this read 
\begin{equation*}
  n_{\mathbf{q}} = \sum_{\mathbf{k}\sigma}\sum_{nn'} \langle n\mathbf{k}\sigma|e^{-i\mathbf{q}\cdot\mathbf{r}}|n'\,\mathbf{k}+\mathbf{q}\,\sigma\rangle \, c^{\dagger}_{n\mathbf{k}\sigma}\, c_{n'\,\mathbf{k}+\mathbf{q}\,\sigma}
\end{equation*}
The many-body definition of this operator is 
\begin{equation*}
  \rho(\mathbf{r}) = \sum_{i=1}^{N} \delta\bigl(\mathbf{r}-\mathbf{r}'_{i}\bigr)
\end{equation*}
Its action can by considering the spatial integral of the function, which is expected to give value only if $\mathbf{r} $ is at particle position $\mathbf{r}'$.
In the momentum eigenstate, the second quantization reads
\begin{equation*}
  \rho (\mathbf{r}) = \frac{1 }{V }\sum_{\mathbf{q }} e ^{i \mathbf{q }\cdot \mathbf{r }}n _{\mathbf{q}} 
\end{equation*}

The transformation is as follows.
\begin{align*}
  \rho(\mathbf{q})
   & = \int d^3r\;e^{-i\mathbf{q}\cdot\mathbf{r}}\;\rho(\mathbf{r})                                                                                                          \\
   & =  \sum_{\mathbf{k},\mathbf{k}'}\left(\int d^3r\;e^{i(\mathbf{k}'-\mathbf{k})\cdot\mathbf{r}}e^{-i\mathbf{q}\cdot\mathbf{r}}\right) \hat{a}_\mathbf{k}^\dagger \hat{a}_{\mathbf{k}'} \\
   & =  \sum_{\mathbf{k},\mathbf{k}'} \delta (\mathbf{k}'-\mathbf{k}-\mathbf{q})   \hat{a}_\mathbf{k}^\dagger \hat{a}_{\mathbf{k}'} \\
   & = \sum_{\mathbf{k}} \hat{a}_\mathbf{k}^\dagger \hat{a}_{\mathbf{k}+\mathbf{q}}= 
   \sum_{\mathbf{k}} \hat{a}_{\mathbf{k}-\mathbf{q}/2}^\dagger \hat{a}_{\mathbf{k}+\mathbf{q}/2}
\end{align*}
Where the last line is obtained by index relabelling $\mathbf{k} \rightarrow \mathbf{k}-\mathbf{q}/2$.

We can do the same without assuming plane wave basis
\begin{equation*}
    \rho(\mathbf{q})
    =\sum_{i}^{N}  \int d^3r\;e^{-i\mathbf{q}\cdot\mathbf{r}}\delta(\mathbf{r }-\mathbf{r}_i) 
    =\sum_{i}^{N}  e^{-i\mathbf{q}\cdot\mathbf{r}_i}
\end{equation*}
Choose the Bloch basis $\ket{n \mathbf{k }\sigma}$, then
\begin{equation*}
  n_{\mathbf{q}}=\sum_{n \mathbf{k} \sigma}\sum_{n' \mathbf{k}k' \sigma'} \langle n \mathbf{k} \sigma|e^{-i\mathbf{q}\cdot\mathbf{r}}|n' \mathbf{k}' \sigma'\rangle c_{n \mathbf{k} \sigma}^{\dagger}c_{n' \mathbf{k}' \sigma'}
\end{equation*}
Use Bloch wavefunctions $\psi_{n\mathbf{k}}(\mathbf{r}) = e^{i\mathbf{k}\cdot\mathbf{r}}\,u_{n\mathbf{k}}(\mathbf{r})$ to compute the matrix elements
\begin{equation*}
 \langle n \mathbf{k} \sigma|e^{-i\mathbf{q}\cdot\mathbf{r}}|n' \mathbf{k}' \sigma'\rangle
  = \delta_{\sigma\sigma'}\int \mathrm{d}^3 r\,u_{n\mathbf{k}}^*(\mathbf{r})\,u_{n'\mathbf{k}'}(\mathbf{r})\,e^{i(\mathbf{k}'-\mathbf{k}-\mathbf{q})\cdot\mathbf{r}}
\end{equation*}
The Bloch function $u_{n\mathbf{k}}(\mathbf{r}+\mathbf{R}) = u_{n\mathbf{k}}(\mathbf{r})$ is periodic, so we can rewrite the integral as 
\begin{equation*}
\sum_{\mathbf R} \int_{\mathrm{cell}} d^3r\; u_{n\mathbf k}^*(\mathbf r+\mathbf R)\, u_{n'\mathbf k'}(\mathbf r+\mathbf R)\, e^{i(\mathbf k'-\mathbf k-\mathbf q)\cdot(\mathbf r+\mathbf R)}
\end{equation*}
The term
\begin{equation*}
 \sum_{\mathbf{R}} 
 e^{i(\mathbf k'-\mathbf k-\mathbf q)\cdot \mathbf R} 
\end{equation*}
produce the selection rule $\mathbf k'-\mathbf k-\mathbf q = \mathbf{G}$.
Restricting to the first Brillouin zone $\mathbf{G}=0$ gives $\mathbf{k}' = \mathbf{k} +\mathbf{q}$.
Under this result, our matrix element reads
\begin{equation*}
     \langle n \mathbf{k} \sigma|e^{-i\mathbf{q}\cdot\mathbf{r}}|n' \mathbf{k}' \sigma'\rangle    
=  \delta_{\sigma\sigma'}\,\delta_{\mathbf{k}',\mathbf{k}+\mathbf{q}}
   \langle n \mathbf{k} \sigma|e^{-i\mathbf{q}\cdot\mathbf{r}}|n' \mathbf{k}+\mathbf{q} \sigma'\rangle    
\end{equation*}
Substituting this matrix element
\begin{equation*}
  n_{\mathbf{q}} = \sum_{\mathbf{k}\sigma}\sum_{nn'} \langle n\mathbf{k}\sigma|e^{-i\mathbf{q}\cdot\mathbf{r}}|n'\,\mathbf{k}+\mathbf{q}\,\sigma\rangle \, c^{\dagger}_{n\mathbf{k}\sigma}\, c_{n'\,\mathbf{k}+\mathbf{q}\,\sigma}
\end{equation*}

Next is the quantization.
The matrix elements reads
\begin{equation*}
  \braket{\mathbf{k }' | \delta(\mathbf{r }-\mathbf{r }) |\mathbf{k}}
  =\frac{1}{V}\int d^3r'\; e^{-i\mathbf{k}'\cdot\mathbf{r}'}\delta(\mathbf{r}-\mathbf{r}')e^{i\mathbf{k}\cdot\mathbf{r}'} = \frac{1 }{V } e ^{i(\mathbf{k }'-\mathbf{k  })\cdot \mathbf{r}}
\end{equation*}
Then  substituting back 
\begin{equation*}
  \rho(\mathbf{r})
  = \frac{1 }{V } \sum_{\mathbf{k},\mathbf{k}'}  e ^{i(\mathbf{k }'-\mathbf{k  })\cdot \mathbf{r}} a ^\dagger _{\mathbf{k }' } a_{\mathbf{k}} 
 =\frac{1 }{V }\sum_{\mathbf{q }} e ^{i \mathbf{q }\cdot \mathbf{r }}n _{\mathbf{q}} 
\end{equation*}

\subsubsection{Example: particle current.}
The particle current, defined as
\begin{equation*}
  \mathbf{j}(\mathbf{r}) = -\frac{i}{2m}\left[\psi^\dagger(\mathbf{r})[\nabla\psi(\mathbf{r})] - [\nabla\psi^\dagger(\mathbf{r})]\;\psi(\mathbf{r})\right]
\end{equation*} 
has the Fourier transform of 
\begin{equation*}
  j(\mathbf{q}) 
= \frac{1}{m}\sum_{\mathbf{k}} \left( \mathbf{k}+\frac{\mathbf{q } }{2} \right)  \hat{a}_{\mathbf{k}}^\dagger \hat{a}_{\mathbf{k}+\mathbf{q}} 
= \frac{1}{m}\sum_{\mathbf{k}} \mathbf{k}\; \hat{a}_{\mathbf{k}-\mathbf{q}/2}^\dagger \hat{a}_{\mathbf{k}+\mathbf{q}/2}
\end{equation*}

The derivation is as follows
\begin{align*}
\mathbf{j}(\mathbf{r}) 
&= -\frac{i}{2m}\sum_{\mathbf{k},\mathbf{k}'}\Big[ \hat{a}_\mathbf{k}^\dagger e^{-i\mathbf{k}\cdot\mathbf{r}}\nabla \hat{a}_{\mathbf{k}'} e^{i\mathbf{k}'\cdot\mathbf{r}}-\big(\nabla  \hat{a}_\mathbf{k}^\dagger e^{-i\mathbf{k}\cdot\mathbf{r}}\big)\cdot \hat{a}_{\mathbf{k}'}e^{i\mathbf{k}'\cdot\mathbf{r}}\Big] \\
&= -\frac{i}{2m}\sum_{\mathbf{k},\mathbf{k}'}\Big[ \hat{a}_\mathbf{k}^\dagger e^{-i\mathbf{k}\cdot\mathbf{r}}\big(i\mathbf{k}'\;a_{\mathbf{k}'}e^{i\mathbf{k}'\mathbf{r}}\big)
-\big(-i\mathbf{k}\;a_\mathbf{k}^\dagger e^{-i\mathbf{k}\cdot\mathbf{r}}\big)\cdot \hat{a}_{\mathbf{k}'}e^{i\mathbf{k}'\cdot\mathbf{r}}\Big] \\
&= \frac{1}{2m}\sum_{\mathbf{k},\mathbf{k}'}(\mathbf{k}+\mathbf{k}')\;a_\mathbf{k}^\dagger \hat{a}_{\mathbf{k}'}\;e^{i(\mathbf{k}'-\mathbf{k})\cdot\mathbf{r}}
\end{align*}
Then apply the transformation
\begin{align*}
j(\mathbf{q}) &= \int d^3r\; e^{-i\mathbf{q}\cdot\mathbf{r}}\; j(\mathbf{r})
= \frac{1}{2m}\sum_{\mathbf{k},\mathbf{k}'} (\mathbf{k}+\mathbf{k}')\; \hat{a}_{\mathbf{k}}^\dagger \hat{a}_{\mathbf{k}'} \int d^3r\; e^{i(\mathbf{k}'-\mathbf{k}-\mathbf{q})\cdot\mathbf{r}}\\
& =  \frac{1}{2m}\sum_{\mathbf{k},\mathbf{k}'} (\mathbf{k}+\mathbf{k}')\; \hat{a}_{\mathbf{k}}^\dagger \hat{a}_{\mathbf{k}'} \delta(\mathbf{k}'-\mathbf{k }-\mathbf{q})\\
&= \frac{1}{2m}\sum_{\mathbf{k}} (\mathbf{k}+(\mathbf{k}+\mathbf{q}))\; \hat{a}_{\mathbf{k}}^\dagger \hat{a}_{\mathbf{k}+\mathbf{q}} 
= \frac{1}{2m}\sum_{\mathbf{k}} (2\mathbf{k}+\mathbf{q})\; \hat{a}_{\mathbf{k}}^\dagger \hat{a}_{\mathbf{k}+\mathbf{q}} \\
j(\mathbf{q}) 
&= \frac{1}{m}\sum_{\mathbf{k}} \left( \mathbf{k}+\frac{\mathbf{q } }{2} \right)  \hat{a}_{\mathbf{k}}^\dagger \hat{a}_{\mathbf{k}+\mathbf{q}} 
= \frac{1}{m}\sum_{\mathbf{k}} \mathbf{k}\; \hat{a}_{\mathbf{k}-\mathbf{q}/2}^\dagger \hat{a}_{\mathbf{k}+\mathbf{q}/2}
\end{align*}
where the last line is obtained by index shift $\mathbf{k}\rightarrow \mathbf{k}-\mathbf{q}/2$.

\subsubsection{Example: position operator.}
Position operator, defined as $\hat{o }^{(1)}=\mathbf{\hat{r}}$, has the one-particle operator form of $\mathbf{R} = \hat{O}_1$.
We evaluated the expression in position basis
\begin{equation*}
  \mathbf{R} = \int d^3r\;d^3r'\; \langle \mathbf{r}|\hat{\mathbf{r}} |\mathbf{r}'\rangle\psi^\dagger(\mathbf{r})\psi(\mathbf{r}')
  = \int d^3r\;\mathbf{r}\;\psi^\dagger(\mathbf{r})\psi(\mathbf{r}) 
\end{equation*}

\subsubsection{Example: linear momentum operator.}
The one-particle operator of linear momentum $P$ has the representation in momentum eigenstate $\ket{\left\{ \mathbf{k} \right\} }$ and position eigenstate $\ket{\left\{ \mathbf{r} \right\} }$ as
\begin{equation*}
  \mathbf{P}
  = \sum_{\mathbf{k }} \mathbf{k}  \hat{a}^\dagger_{\mathbf{k }} \hat{a}_{\mathbf{k }}    
  =-i \int d^{3}r\; \psi^{\dagger}(\mathbf{r})[\nabla\psi(\mathbf{r})]
\end{equation*}

The derivation in momentum eigenstate is
\begin{equation*}
  \mathbf{P} = \sum_{\mathbf{k },\mathbf{k}'} \braket{\mathbf{k } | \mathbf{p }|\mathbf{k}'} \hat{a}^\dagger_{\mathbf{k }} \hat{a}_{\mathbf{k }'}    
  = \sum_{\mathbf{k },\mathbf{k}'} \mathbf{k}' \delta_{\mathbf{k}\mathbf{k}'} \hat{a}^\dagger_{\mathbf{k }} \hat{a}_{\mathbf{k }'}    
  = \sum_{\mathbf{k }} \mathbf{k}  \hat{a}^\dagger_{\mathbf{k }} \hat{a}_{\mathbf{k }}    
\end{equation*}
while in position eigenstate
\begin{equation*}
  \mathbf{P}
  = \int d^3r\;\psi^\dagger(\mathbf{r}) \langle \mathbf{r}| \mathbf{p} |\mathbf{r}\rangle \psi(\mathbf{r})
 =-i   \int d^3r\;d^3r'\;\psi^\dagger(\mathbf{r})\delta(\mathbf{r}-\mathbf{r'}) [\nabla   \psi(\mathbf{r}')]
\end{equation*}
and the Dirac function collapse the primed integral.

\subsubsection{Example: one-particle Hamiltonian.}
Consider the one–particle Hamiltonian
\begin{equation*}
  \hat{h} = \frac{\hat{\mathbf{p}}^{2}}{2m} + V(\mathbf{r})
\end{equation*}
In the momentum eigenbasis, the one-particle operator reads
\begin{equation*}
  H = \sum_{\mathbf{k}} \frac{\mathbf{k}^{2}}{2m} a_{\mathbf{k}}^{\dagger} a_{\mathbf{k}}    + \sum_{\mathbf{k} \mathbf{k}'} V_{\mathbf{q}} a_{\mathbf{k} + \mathbf{q}}^{\dagger} a_{\mathbf{k}}
\end{equation*}
with $V_{\mathbf{q}}$ as Fourier component of the spatial potential that scatters a particle from momentum $\mathbf{k}'$ to $\mathbf{k}$ with momentum transfer $\mathbf{q}=\mathbf{k}-\mathbf{k}'$
\begin{equation*}
  V_\mathbf{q} = \frac{1 }{V } \int d^3r\; e^{- i\mathbf{q }\cdot \mathbf{r}} V(\mathbf{r})
\end{equation*}
From the action of annihilation and creation operator, it can also be interpreted that $a_{\mathbf{k}} $ annihilate particle with momentum $\mathbf{k }$, then particle with momentum $\mathbf{k}+\mathbf{q}$ is created by $a_{\mathbf{k }+\mathbf{q}}$.
In the position eigenbasis
\begin{equation*}
  H = \int d^3r\; \psi^\dagger(\mathbf r)\left(-\frac{1}{2m}\nabla^2 + V(\mathbf r)\right)\psi(\mathbf r)
\end{equation*}

The derivation in position eigenstate is quite straightforward since the Hamiltonian is local in position representation.
For the momentum eigenstate, the kinetic term  reads  
\begin{equation*}
    K=\frac{1 }{2m} \sum_{\mathbf{k},\mathbf{k}'} \langle k|\hat{\mathbf{p}}^2 |k'\rangle a^\dagger_{\mathbf{k}} a_{\mathbf{k}'} = \sum_{\mathbf{k}} \frac{k^2}{2m} a^\dagger_{\mathbf{k}} a_{\mathbf{k}}
\end{equation*}
For the potential term 
\begin{equation*}
  V = \sum_{\mathbf{k},\mathbf{k}'} \langle \mathbf{k}|V(\mathbf{r})|\mathbf{k}'\rangle  a_{\mathbf{k}}^\dagger a_{\mathbf{k}'}
\end{equation*}
We evaluate the element in position representation, since the wavefunction is known to be the plane wave 
\begin{align*}
    \langle \mathbf{k}|V(\mathbf{r})|\mathbf{k}'\rangle &= \frac{1}{V}\int d^3r\;e^{-i(\mathbf{k}-\mathbf{k}')\cdot\mathbf{r}}V(\mathbf{r})\\
    V_{\mathbf{k} - \mathbf{k}'} = V_{\mathbf{q}} &= \frac{1}{V}\int d^3r\;e^{-i\mathbf{q}\cdot\mathbf{r}}V(\mathbf{r})
\end{align*}
This is just Fourier transform of the potential.

\subsubsection{Example: Spin operator.}
For a system of $N$ electrons
\begin{equation*}
  \mathbf{S} = \sum_{i=1}^{N} \mathbf{S}_i
\end{equation*}
where $\mathbf{S }$ act on particle $i$.
Let $\mathbf{S}=(S_x,S_y,S_z)$, then the second quantization version is 
\begin{equation*}
  \mathbf{S} = \frac{\hbar}{2}\sum_{k}\Bigl(\bigl[c_{\mathbf{k}\uparrow}^\dagger c_{\mathbf{k}\downarrow}+c_{\mathbf{k}\downarrow}^\dagger c_{\mathbf{k}\uparrow}\,\bigr],
i\bigl[c_{\mathbf{k}\downarrow}^\dagger c_{\mathbf{k}\uparrow}-c_{\mathbf{k}\uparrow}^\dagger c_{\mathbf{k}\downarrow}\bigr], 
\bigl[c_{\mathbf{k}\uparrow}^\dagger c_{\mathbf{k}\uparrow}-c_{\mathbf{k}\downarrow}^\dagger c_{\mathbf{k}\downarrow}\bigr]\Bigr)
\end{equation*}
More generally
\begin{equation*}
  \mathbf{\hat{S}} = \frac{\hbar}{2}\sum_{\mathbf{k}}\sum_{\sigma,\sigma'} a^\dagger_{\mathbf{k}\sigma}\boldsymbol{\sigma}_{\sigma\sigma'}a_{\mathbf{k}\sigma'}
\end{equation*}
where $\boldsymbol{\sigma}_{\sigma\sigma'} = \braket{\sigma  | \boldsymbol{\sigma }|\sigma'}$ is the matrix element in $\ket{\mathbf{k }\sigma}$ basis.

The derivation start with recognizing $\mathbf{S }= \hbar \boldsymbol{\sigma }/2$.
Insert into the general one-body operator formula to obtain the general form
\begin{equation*}
  \mathbf{\hat{S} } 
  = \sum_{\mathbf{k}\sigma} \sum_{\mathbf{k}'\sigma'}  \langle \mathbf{k}\sigma|\mathbf{\hat{S}}|\mathbf{k}'\sigma'\rangle a^\dagger_{\mathbf{k}\sigma} a_{\mathbf{k}'\sigma'} 
  = \sum_{\mathbf{k}}\sum_{\sigma,\sigma'} \frac{\hbar}{2}\boldsymbol{\sigma}_{\sigma\sigma'} a^\dagger_{\mathbf{k}\sigma} a_{\mathbf{k}\sigma'}
\end{equation*}
The explicit component is evaluated using these components of Pauli matrices
\begin{equation*}
    \sigma_1 =
    \begin{bmatrix}
        0 & 1 \\
        1 & 0
    \end{bmatrix}
    \quad
    \sigma_2 =
    \begin{bmatrix}
        0 & -i \\
        i & 0
    \end{bmatrix}
    \quad
    \sigma_3    =
    \begin{bmatrix}
        1 & 0  \\
        0 & -1
    \end{bmatrix}
\end{equation*}a
In spin $1/2$, all elements is evaluated as 
\begin{equation*}
  \sigma_n = 
\begin{bmatrix}
\langle\uparrow|\sigma_n|\uparrow\rangle & \langle\uparrow|\sigma_n|\downarrow\rangle\\[6pt]
\langle\downarrow|\sigma_n|\uparrow\rangle & \langle\downarrow|\sigma_n|\downarrow\rangle
\end{bmatrix}.
\end{equation*}
and the rest is trivial.

\subsubsection{Derivation: second quantization representation.}
First consider the action of $\hat{o} ^{1}$ on one-particle basis $\ket{\left\{ 1_\alpha \right\} }$, which we assume as its eigenstate
\begin{equation*}
  \hat{o } ^{(1)}\ket{1_\alpha } =  \omega_\alpha \ket{1_\alpha}
\end{equation*}
Recall that the one-particle operator $\hat{O }_1$ is defined as the sum of this operator.
Hence, its action on this basis is
\begin{equation*}
  \hat{O }_1 =  \sum_{\alpha } \omega_\alpha n_\alpha = \sum_{\alpha } \omega_\alpha \hat{a}_\alpha ^\dagger \hat{a}_\alpha
\end{equation*}
The one-particle operator is also diagonal in its basis, hence we can write the eigenvalue as the matrix element
\begin{equation*}
  \hat{O }_1
  = \sum_{\alpha } \braket{1_\alpha  | \hat{o } ^{(1)}| 1_\alpha} \hat{a}_\alpha ^\dagger \hat{a}_\alpha
\end{equation*}
Transform into arbitrary one-particle basis by using the completeness relations
\begin{multline*}
    \hat{o }_1
  =\sum_{i,j}  \sum_{\alpha } \braket{1_\alpha  | 1_i} \braket{1_i | \hat{o } ^{(1)}|1_j   } \braket{1_j  | 1_\alpha} \hat{a}_\alpha ^\dagger \hat{a}_\alpha\\
  =\sum_{i,j}  \sum_{\alpha } \braket{1_i | \hat{o } ^{(1)}|1_j   } \delta_{\alpha i} \hat{a}_\alpha ^\dagger \delta_{\alpha j}a_\alpha
  =\sum_{i,j}  \ \braket{1_i | \hat{o } ^{(1)}|1_j   }  \hat{a}_i ^\dagger a_j
\end{multline*}
\begin{equation*}
\end{equation*}
The $\alpha$ index collapse, then the representation is obtained.

\subsection{Two-particle operator}
A general two-particle operator is
\begin{equation*}
  \hat{O }_2 = \sum_{ij } \hat{o }_{ij }^{(2)}
\end{equation*}
where the sum runs over all pairs of particles.
The second quantized act in the abstract two-particle basis
\begin{equation*}
  \hat{O}_2 = \sum_{i,j,m,n} \braket{1_i 1_j  | \hat{o} ^{(2)}| 1_m 1_n } \hat{a}^\dagger_{i}  \hat{a}^\dagger_{j} \hat{a}_n \hat{a}_m
\end{equation*}
Note the reversed order of the indices $n$ and $m$ of the pair of annihilation operators as compared to their order in the matrix elements, which we adopt that due to Fermions antisymmetry.
This is imposed to enforce that the operator ordering $a_na_m$ exactly removes the particles in states $m$ and $n$ without producing an additional minus sign.

In the position basis, the expression reads
\begin{equation*}
  \hat{O}_2 = \int d^3 \mathbf{r}_1 \; d^3 \mathbf{r}_1' \; d^3 \mathbf{r}_2 \; d^3 \mathbf{r}_2' \; \psi^\dagger(\mathbf{r}_1) \psi(\mathbf{r}_1') \langle \mathbf{r}_1, \mathbf{r}_1' | \hat{o}^{(2)} | \mathbf{r}_2, \mathbf{r}_2' \rangle \psi(\mathbf{r}_2') \psi(\mathbf{r}_2)
\end{equation*}
If $\hat{O }_2$ is diagonal
\begin{equation*}
  \langle \mathbf{r}_1, \mathbf{r}_1' | \hat{O}^{(2)} | \mathbf{r}_2, \mathbf{r}_2' \rangle = \hat{O}^{(2)}(\mathbf{r}_1, \mathbf{r}_1') \delta(\mathbf{r}_1 - \mathbf{r}_2) \delta(\mathbf{r}_1' - \mathbf{r}_2')
\end{equation*}
the expression reduces further
\begin{equation*}
  \hat{O}_2 = \int d^3 \mathbf{r} d^3 \mathbf{r}' \; \psi^\dagger(\mathbf{r}) \psi(\mathbf{r}') \hat{o}^{(2)}(\mathbf{r}, \mathbf{r}') \psi(\mathbf{r}') \psi(\mathbf{r})
\end{equation*}

Imposing translational invariance $\hat{o}^{(2)}(\mathbf{r}, \mathbf{r}') = \hat{o}^{(2)}(\mathbf{r} - \mathbf{r}')$, we obtain
\begin{align*}
  \hat{O}_2
   & =  \sum_{\mathbf{k}_1 \mathbf{k}_2 } \sum_{\mathbf{q}} \hat{o}^{(2)}(\mathbf{q}) \hat{a}_{\mathbf{k}_1}^\dagger \hat{a}_{\mathbf{k}_2}^\dagger \hat{a}_{\mathbf{k}_2+\mathbf{q}} \hat{a}_{\mathbf{k}_1-\mathbf{q}} \\
  \hat{O}_2
   & = \sum_{\mathbf{k}_1 \mathbf{k}_2 } \sum_{\mathbf{q}} \hat{o}^{(2)}(\mathbf{q}) \hat{a}_{\mathbf{k}_1+\mathbf{q}}^\dagger \hat{a}_{\mathbf{k}_2-\mathbf{q}}^\dagger \hat{a}_{\mathbf{k}_2} \hat{a}_{\mathbf{k}_1}
\end{align*}
where \( \hat{o}^{(2)}(\mathbf{q}) \) is the Fourier transform of \( \hat{o}^{(2)}(\mathbf{r}) \)
\begin{equation*}
  \hat{o}^{(2)}(\mathbf{q}) = \int d^3 \mathbf{r}\; \hat{o}^{(2)}(\mathbf{r}) e^{-i\mathbf{q} \cdot r}
  \qquad
  \begin{cases}
    \mathbf{r} = \mathbf{r}_1-\mathbf{r}_2\\
    \mathbf{q} = \mathbf{k}_1 + \mathbf{k}_1' \\
    \mathbf{q}= \mathbf{k}_2' - \mathbf{k}_2
  \end{cases}
\end{equation*}

\subsubsection{State order reversal.}
As stated before, the purpose of the reversed ordering is to ensure that the annihilation operators remove the particles in the ket in the same order in which they were created, thereby avoiding an additional fermionic sign.
To see this, consider the correctly ordered operators
\begin{align*}
a_n a_m \lvert 1_m 1_n\rangle 
&= a_n a_m a_m^\dagger a_n^\dagger \lvert 0\rangle,\\
& =  a_n(1 - a_m^\dagger a_m) a_n^\dagger \lvert 0\rangle \\
& =  (a_n a_n^\dagger -  a_m^\dagger a_m a_n^\dagger )  \lvert 0\rangle \\
& =  a_n a_n^\dagger   \lvert 0\rangle + a_m^\dagger a_n^\dagger a_m  \lvert 0\rangle \\
a_n a_m \lvert 1_m 1_n\rangle 
&= \lvert 0\rangle
\end{align*}
If the order was reversed, however
\begin{align*}
a_m a_n|1_m1_n\rangle 
&= a_m a_n a_m^\dagger a_n^\dagger |0\rangle\\
&= -a_m a_m^\dagger a_n a_n^\dagger|0\rangle\\
&= -\big(1-a_m^\dagger a_m\big)a_n a_n^\dagger|0\rangle\\
&= - a_n a_n^\dagger|0\rangle + a_m^\dagger a_m a_n a_n^\dagger|0\rangle\\
a_m a_n|1_m1_n\rangle&= -|0\rangle
\end{align*}
an extra negative sign was produced

\subsubsection{Example: Coulomb potential.}
Two-particle Coulomb interaction is 
\begin{equation*}
  \hat{o}^{(2)} 
  = \frac{e^2}{4\pi\varepsilon_0}\frac{1}{|{\mathbf r}_1-{\mathbf r}_2|} 
\end{equation*}
Then the second quantization of Coulomb potential is 
\begin{equation*}
  \hat{V}
  = \frac{1}{2}\sum_{i,j,m,n}V_{ijmn} a_i^\dagger a_j^\dagger a_n a_m 
\end{equation*}
with $u_i(\mathbf{r}_n) = \braket{\mathbf{r}_n  | 1_i}$ as the single particle orbital and factor $1/2$ to avoid double counting.
The derivation is quite straightforward, with the integral term 
\begin{equation*}
  V_{ijmn}
  = \int d^3r_1\;d^3r_2 u_i^*({\mathbf r}_1)u_j^*({\mathbf r}_2)\frac{e^2}{4\pi\varepsilon_0|{\mathbf r}_1-{\mathbf r}_2|}u_m({\mathbf r}_1)u_n({\mathbf r}_2)
\end{equation*}
as the matrix element of the Coulomb potential in that basis.

Writing in the momentum space, yields
\begin{align*}
  \hat{V } 
  & = \frac{1}{2V} \sum_{\mathbf{k}_1,\mathbf{k}_2,\mathbf{q}} V(\mathbf{q}) a_{\mathbf{k}_1+\mathbf{q}}^\dagger a_{\mathbf{k}_2-\mathbf{q}}^\dagger a_{\mathbf{k}_2} a_{\mathbf{k}_1}\\
  \hat{V } 
  & = \frac{1}{2V} \sum_{\mathbf{k}_1,\mathbf{k}_2,\mathbf{q}} \frac{e^2 }{\epsilon_0 q^2} a_{\mathbf{k}_1+\mathbf{q}}^\dagger a_{\mathbf{k}_2-\mathbf{q}}^\dagger a_{\mathbf{k}_2} a_{\mathbf{k}_1}
\end{align*}
\begin{equation*}
\end{equation*}
with $V(\mathbf{q})$ as the Fourier transform of the spatial Coulomb potential
\begin{equation*}
V(\mathbf{q }) 
= \int d^3r\;    \frac{e^2}{4\pi\varepsilon_0}\frac{1}{|{\mathbf r}_1-{\mathbf r}_2|}   
=\frac{e^2}{4\pi\varepsilon_0} \int d^3r\;    \frac{e ^{i \mathbf{q }\cdot \mathbf{r}}}{r}   
= \frac{e^2 }{\epsilon_0 q^2}
\end{equation*}

The integral evaluation is as follows
\begin{align*}
    I(q) & =  \int d^3 r \; \frac{e^{i \mathbf{q }\cdot \mathbf{ r }}}{r } \\
  & = \int_{0}^{\infty}dr\;r^{2}\int_{0}^{\pi}d\theta\;\sin\theta\int_{0}^{2\pi}d\phi\;\frac{1}{r}e^{iqr\cos\theta} \\
  I(q) & =  2\pi\int_{0}^{\infty}dr\;r\int_{0}^{\pi}d\theta\sin\theta e^{iqr\cos\theta}
\end{align*}
The second term can be evaluated by substituting $u = \cos \theta$ and $du = -\sin \theta\;d\theta$
\begin{equation*}
\int_{0}^{\pi} d\theta\;\sin\theta\;e^{iqr\cos\theta}
= \int_{-1}^{1} du\; e^{iqru}
= \frac{e^{iqr}-e^{-iqr}}{iqr}
= \frac{2\sin(qr)}{qr}
\end{equation*}
The integral now reads
\begin{equation*}
  I(q) 
  = \frac{4\pi}{q}\int_{0}^{\infty}dr\;\sin(qr)
\end{equation*}
This can be evaluated using Feynman's method of parameterization 
\begin{equation*}
  I(\epsilon) 
  = \int_{0 }^{\infty}dr\; e ^{(iq -\epsilon )r}
\end{equation*}
with real $e ^{-\epsilon r}$. 
Then, to evaluate the integral, we simply have to 
\begin{equation*}
  \int_{0}^{\infty}dr\;\sin(qr)  
  = \lim_{\epsilon \rightarrow 0} \mathrm{Im} \left(\int_{0 }^{\infty}dr\; e ^{(iq -\epsilon )r}\right) 
\end{equation*}
Evaluate the function
\begin{equation*}
  I(\epsilon)  
  =\left[-\frac{1}{\epsilon-iq}e^{-(\epsilon-iq)r}\right]_{0}^{\infty}
  =\frac{1}{\epsilon-iq}
  = \frac{\epsilon+iq }{\epsilon^2 + q^2}
\end{equation*}
Taking the imaginary component we obtain the integral value 
\begin{equation*}
  I(q) \frac{2\pi }{q }\frac{1 }{q } = \frac{4\pi }{q^2}
\end{equation*}
Multiplying this with the other Coulomb constant, we obtain the Fourier transform of the Coulomb potential. 


\subsubsection{Derivation: momentum representation.}
Translational invariance requires that shifting the entire system by a vector $\mathbf{a}$ leaves the interaction unchanged
\begin{equation*}
  V(\mathbf{r}_1+\mathbf{a},\mathbf{r}_2+\mathbf{a})=
  V(\mathbf{r}_1,\mathbf{r}_2)
\end{equation*}
The only functions satisfying this condition depend solely on the relative coordinate
\begin{equation*}
  V(\mathbf{r}_1,\mathbf{r}_2)
  =   V(\mathbf{r}_1-\mathbf{r}_2)
\end{equation*}
If the interaction is a pair potential depending only on the separation
\begin{equation*}
  \langle \mathbf{r}_1,\mathbf{r}_2 | \hat{o}^{(2)} | \mathbf{r}_1', \mathbf{r}_2' \rangle
  = \hat{o}^{(2)}(\mathbf{r}_1 - \mathbf{r}_2) \delta(\mathbf{r}_1 - \mathbf{r}_1') \delta(\mathbf{r}_2 - \mathbf{r}_2')
\end{equation*}
First we start with
\begin{equation*}
  \hat{O}_2 = \sum_{\mathbf{k}_1, \mathbf{k}_2} \sum_{\mathbf{k}_1' \mathbf{k}_2'}  \langle \mathbf{k}_1 \mathbf{k}_2 | \hat{o}^{(2)} | \mathbf{k}_2' \mathbf{k}_1' \rangle a ^\dagger _{\mathbf{k}_1} a ^\dagger _{\mathbf{k}_2 } a _{\mathbf{k}_2'} a _{\mathbf{k}_1'}
\end{equation*}
The element can be evaluated as
\begin{multline*}
  \langle \mathbf{k}_1 \mathbf{k}_2 | \hat{o}^{(2)} | \mathbf{k}_1' \mathbf{k}_2' \rangle \\
  = \int d^3 \mathbf{r}_1\; d^3 \mathbf{r}_2\; d^3 \mathbf{r}_{1}'\; d^3 \mathbf{r}_{2}'\; \langle \mathbf{k}_1 \mathbf{k}_2 | \mathbf{r}_1, \mathbf{r}_2 \rangle \langle \mathbf{r}_1, \mathbf{r}_2 | \hat{o}^{(2)} | \mathbf{r}_1', \mathbf{r}_2'\rangle\langle \mathbf{r}_1',\mathbf{r}_2' | \mathbf{k}_1' \mathbf{k}_2' \rangle
\end{multline*}
Then the integral reduces
\begin{multline*}
  \langle \mathbf{k}_1 \mathbf{k}_2 | \hat{o}^{(2)} | \mathbf{k}_1' \mathbf{k}_2' \rangle \\
  = \int d^3 \mathbf{r}_1\; d^3 \mathbf{r}_2\; \langle \mathbf{k}_1 \mathbf{k}_2 | \mathbf{r}_1, \mathbf{r}_2 \rangle \langle \mathbf{r}_1, \mathbf{r}_2 | \hat{o}^{(2)} | \mathbf{r}_1, \mathbf{r}_2\rangle\langle \mathbf{r}_1,\mathbf{r}_2 | \mathbf{k}_1' \mathbf{k}_2' \rangle
\end{multline*}
In this composite state,  the plane wave representation reads as
\begin{equation*}
  \langle \mathbf{r}_1,\mathbf{r}_2 | \mathbf{k}_1 \mathbf{k}_2 \rangle
  = \frac{1 }{\sqrt{V}}  e ^{i(\mathbf{k}_1 \cdot  \mathbf{r}_1+\mathbf{k}_2 \cdot  \mathbf{r}_2)}
\end{equation*}
Substituting back
\begin{equation*}
  \frac{1}{V} \int d^3 \mathbf{r}_1 d^3 \mathbf{r}_2 \hat{o}^{(2)}(\mathbf{r}_1 -\mathbf{r}_2) e^{-i(\mathbf{k}_1 - \mathbf{k}_1') \cdot \mathbf{r}_1} e^{-i(\mathbf{k}_2 -\mathbf{k}_2') \cdot \mathbf{r}_2}
\end{equation*}
Introduce new coordinates
\begin{equation*}
  \begin{cases}
    \mathbf{R }=\mathbf{r }_1 +\mathbf{r}_2 \\
    \mathbf{r }=\mathbf{r }_1 -\mathbf{r}_2
  \end{cases}
  \qquad
  \text{or}
  \qquad
  \begin{cases}
    \mathbf{r}_1=\mathbf{R } +\mathbf{r} \\
    \mathbf{r}_2=\mathbf{R } -\mathbf{r}
  \end{cases}
\end{equation*}
Rewrite the exponent as
\begin{align*}
  (\mathbf{k}_1 - \mathbf{k}_1') \cdot \mathbf{r}_1 + (\mathbf{k}_2 - \mathbf{k}_2') \cdot \mathbf{r}_2
   & =  (\mathbf{k}_1 + \mathbf{k}_2 - \mathbf{k}_1' - \mathbf{k}_2') \cdot \mathbf{R }                     \\
   & \qquad+ \left[ (\mathbf{k}_1 - \mathbf{k}_1') - (\mathbf{k}_2 - \mathbf{k}_2') \right]\cdot \mathbf{r} \\
  (\mathbf{k}_1 - \mathbf{k}_1') \cdot \mathbf{r}_1 + (\mathbf{k}_2 - \mathbf{k}_2') \cdot \mathbf{r}_2
   & = \mathbf{p }\cdot \mathbf{R }+ \mathbf{q }\cdot \mathbf{r}
\end{align*}
such that the integral separates
\begin{align*}
  \langle \mathbf{k}_1 \mathbf{k}_2 | \hat{o}^{(2)} | \mathbf{k}_1' \mathbf{k}_2' \rangle
   & = \frac{1 }{V}  \int d^3 R \; e^{-i \mathbf{p}\cdot \mathbf{R}} \int d^3 \mathbf{r}\; \hat{o}^{(2)}(\mathbf{r}) e^{-i \mathbf{q} \cdot \mathbf{r}} \\
  \langle \mathbf{k}_1 \mathbf{k}_2 | \hat{o}^{(2)} | \mathbf{k}_1' \mathbf{k}_2' \rangle
   & =  \delta (\mathbf{p}) \int d^3 \mathbf{r}\; \hat{o}^{(2)}(\mathbf{r}) e^{-i \mathbf{q} \cdot \mathbf{r}}
\end{align*}
The delta function enforce $\mathbf{k }_1+\mathbf{k }_2=\mathbf{k }_1'+\mathbf{k }_2'$, so the definition of $\mathbf{q}$ simplifies and the integral reads
\begin{equation*}
  \langle \mathbf{k}_1 \mathbf{k}_2 | \hat{o}^{(2)} | \mathbf{k}_1' \mathbf{k}_2' \rangle
  =
  \int d^3 \mathbf{r}\; \hat{o}^{(2)}(\mathbf{r}) e^{-i \mathbf{q} \cdot \mathbf{r}}
  \qquad
  \begin{cases}
    \mathbf{q} = \mathbf{k}_1 +\mathbf{k}_1' \\
    \mathbf{q}= \mathbf{k}_2' - \mathbf{k}_2
  \end{cases}
\end{equation*}

\end{document}